% ============================================================
%  sections/analysis.tex — Interpretation & discussion
% ============================================================

\begin{tcolorbox}[notebox, title={\faChartLine\quad Analysis \& Interpretation}]

  \textbf{Conductivity vs.\ baseline}\\[4pt]
  {\small
    The measured total ionic conductivity of P-047-A ($\sigma = 1.85 \times 10^{-3}$\,S\,cm\textsuperscript{-1})
    represents a \textbf{16.8-fold improvement} over undoped batch NRI-2024-0041
    ($\sigma = 1.10 \times 10^{-4}$\,S\,cm\textsuperscript{-1}) and surpasses
    the $3.0 \times 10^{-4}$\,S\,cm\textsuperscript{-1} hypothesis threshold by a factor of 6.2.
    This result aligns with literature values for Al-doped cubic LLZO
    \cite{Geiger2011, Murugan2007} and confirms efficient phase stabilisation.
  }

  \vspace{8pt}
  \textbf{Grain boundary resistance}\\[4pt]
  {\small
    The grain-boundary contribution ($R_{\text{gb}} = 312.7\,\Omega$) accounts for
    67.8\,\% of total impedance, indicating grain-boundary resistance as the dominant
    limitation. This is higher than the 55\,\% reported in \cite{Geiger2011} for
    similar compositions, suggesting further optimisation of sintering temperature or
    dwell time may be beneficial. A companion experiment at 1175\,\textdegree C is
    planned for batch NRI-2024-0050.
  }

  \vspace{8pt}
  \textbf{Relative density shortfall}\\[4pt]
  {\small
    Relative density of 92.8\,\% falls below the 94\,\% target. Possible causes:
    \begin{itemize}[leftmargin=16pt, itemsep=2pt, topsep=2pt]
      \item Use of furnace NRI-F-02 (deviation) — thermocouple calibrated but
            temperature profile across sample zone may differ from NRI-F-01.
      \item Insufficient milling homogeneity following jar replacement in Step~3
            (milling jar changeover interrupted the programmed milling cycle by
            $\sim$18\,min).
    \end{itemize}
    Despite the density shortfall, ionic conductivity substantially exceeds the
    target, suggesting that grain-boundary chemistry (Al segregation) dominates
    over microstructural porosity in this composition window.
  }

  \vspace{8pt}

  \noindent
  \begin{minipage}[t]{0.48\linewidth}
    \begin{tcolorbox}[
      enhanced,
      colback=PageTint,
      colframe=AccentGold,
      leftrule=3pt,
      rightrule=0.3pt,
      toprule=0.3pt,
      bottomrule=0.3pt,
      colbacktitle=AccentGold!20,
      coltitle=InkBlue,
      fonttitle=\bfseries\footnotesize,
      title={Key Findings},
      arc=2pt,
      left=8pt, right=8pt,
      top=6pt, bottom=6pt,
    ]
      {\small
        \begin{itemize}[leftmargin=12pt, itemsep=2pt, topsep=0pt]
          \item[\faCheck] Cubic LLZO phase confirmed.
          \item[\faCheck] $\sigma = 1.85\times10^{-3}$\,S\,cm\textsuperscript{-1}.
          \item[\faCheck] Hypothesis exceeded by $\times$6.2.
          \item[\faTimes] Relative density below target (92.8\,\%).
          \item[\faTimes] Pellet P-047-B cracked during pressing.
        \end{itemize}
      }
    \end{tcolorbox}
  \end{minipage}%
  \hfill
  \begin{minipage}[t]{0.48\linewidth}
    \begin{tcolorbox}[
      enhanced,
      colback=PageTint,
      colframe=InkBlue,
      leftrule=3pt,
      rightrule=0.3pt,
      toprule=0.3pt,
      bottomrule=0.3pt,
      colbacktitle=LightSteel,
      coltitle=InkBlue,
      fonttitle=\bfseries\footnotesize,
      title={Next Steps},
      arc=2pt,
      left=8pt, right=8pt,
      top=6pt, bottom=6pt,
    ]
      {\small
        \begin{itemize}[leftmargin=12pt, itemsep=2pt, topsep=0pt]
          \item Rietveld refinement (P.\ Nair, 16 Nov).
          \item 0.4\,mol\% Al batch (NRI-2024-0048).
          \item SEM cross-section of P-047-A.
          \item Repeat with NRI-F-01 once available.
          \item Update project tracking in Synapse LIMS.
        \end{itemize}
      }
    \end{tcolorbox}
  \end{minipage}

\end{tcolorbox}

% ── References ─────────────────────────────────────────────
\vspace{4pt}
\begin{tcolorbox}[sideruled, title={\faBook\quad References}]
  \renewcommand{\refname}{}
  \begingroup
    \small
    \bibliographystyle{unsrt}
    \bibliography{references}
  \endgroup
\end{tcolorbox}
